\documentclass[a4paper,11pt]{article}
\usepackage[utf8]{inputenc}
\usepackage[french]{babel}
\usepackage[T1]{fontenc}
\usepackage{amsmath,amssymb,amsthm}
\usepackage{geometry}
\geometry{margin=2cm}
\usepackage{enumitem}
\usepackage{hyperref}
\usepackage{booktabs}
\usepackage{array}
\usepackage{xcolor}
\usepackage{listings}
\usepackage{titlesec}
\usepackage{fancyhdr}
\usepackage{lastpage}
\usepackage{graphicx}
\usepackage{etoolbox}
\AtBeginDocument{\shorthandoff{;:!?}}

\definecolor{codebg}{gray}{0.95}
\definecolor{codeframe}{gray}{0.80}

\titleformat{\section}{\large\bfseries\color{blue!60!black}}{}{0pt}{}[\vspace{-0.5ex}\rule{\textwidth}{0.4pt}]
\titleformat{\subsection}{\bfseries\color{blue!40!black}}{}{0pt}{}
\titleformat{\subsubsection}{\itshape\color{blue!30!black}}{}{0pt}{}

\usepackage{listings}
\lstdefinestyle{rmono}{
  frame=single,
  numbers=none,
  breaklines=true,
  basicstyle=\small\ttfamily,
  backgroundcolor=\color{codebg},
  rulecolor=\color{codeframe},
  columns=fullflexible,
  keepspaces=true,
  showstringspaces=false,
  language=R
}
\lstset{style=rmono}

\pagestyle{fancy}
\fancyhf{}
\fancyhead[L]{\small Analyse critique --- ROSE (Lunardon, Menardi, Torelli, 2014)}
\fancyhead[R]{\small Fiche r\'ecapitulative}
\fancyfoot[C]{\thepage/\pageref{LastPage}}
\renewcommand{\headrulewidth}{0.4pt}

\theoremstyle{definition}
\newtheorem*{definition}{D\'efinition}
\newtheorem*{remark}{Remarque}
\newtheorem*{important}{\`A retenir}

\title{\textbf{Analyse critique -- ROSE: A Package for Binary Imbalanced Learning}\\[4pt]
\large Lunardon, Menardi, Torelli (R Journal, 6(1), 2014)\\[2pt]
\small Package R pour l'apprentissage sur donn\'ees d\'es\'equilibr\'ees}
\author{\'Etudiant STA211}
\date{14 juin 2026}

\begin{document}
\maketitle
\thispagestyle{empty}

\begin{abstract}
Cet article pr\'esente le package R \textbf{ROSE} (Random Over-Sampling
Examples), d\'edi\'e \`a la classification binaire en pr\'esence de
classes d\'es\'equilibr\'ees. La m\'ethode g\'en\`ere des
\'echantillons artificiels \'equilibr\'es par bootstrap liss\'e,
facilitant \`a la fois l'estimation du mod\`ele et l'\'evaluation de
sa performance. Cette fiche propose une analyse critique structur\'ee,
avec illustrations num\'eriques, extraits de code R et
QCM.
\end{abstract}

\vfill
\tableofcontents
\newpage

% ===================================================================
\section{Pr\'esentation g\'en\'erale}
% ===================================================================

\subsection{Contexte}

Les donn\'ees d\'es\'equilibr\'ees (ex.~2\% de positifs, 98\% de
n\'egatifs) posent un probl\`eme majeur en classification supervis\'ee.
Un classifieur na\"if qui pr\'edit toujours la classe majoritaire
atteint 98\% de bonne classification, mais ne d\'etecte aucun cas rare.

L'article propose le package \textbf{ROSE} qui impl\'emente une
m\'ethode unifi\'ee pour~:
\begin{itemize}
  \item \textbf{Estimer} un classifieur malgr\'e le d\'es\'equilibre
    (via des donn\'ees artificielles \'equilibr\'ees)
  \item \textbf{\'Evaluer} la performance de fa\c con fiable
    (via des estimateurs sp\'ecifiques : holdout, bootstrap,
    validation crois\'ee)
\end{itemize}

\subsection{R\'ef\'erence}

\begin{itemize}
  \item Lunardon N., Menardi G., Torelli N. (2014).
    \emph{ROSE: A Package for Binary Imbalanced Learning}.
    The R Journal, 6(1), 79--89.
  \item M\'ethode sous-jacente~: Menardi G., Torelli N. (2014).
    \emph{Training and assessing classification rules with imbalanced
    data}. Data Mining and Knowledge Discovery, 28(1), 92--122.
\end{itemize}

% ===================================================================
\section{Probl\`eme du d\'es\'equilibre de classes}
% ===================================================================

\subsection{Cons\'equences}

\begin{enumerate}
  \item \textbf{Estimation}~: les r\`egles de classification sont
    \guillemotleft submerg\'ees \guillemotright par la classe
    majoritaire~; la classe rare est ignor\'ee.
  \item \textbf{\'Evaluation}~: le taux d'erreur global (accuracy)
    est trompeur~; on utilise plut\^ot pr\'ecision, rappel,
    $F$-mesure, ou AUC.
  \item \textbf{Seuil}~: le seuil par d\'efaut ($\hat{p} > 0,5$)
    peut ne jamais \^etre d\'epass\'e~; la courbe ROC \'evite ce
    choix arbitraire.
\end{enumerate}

\subsection{Rem\`edes traditionnels}

\begin{itemize}
  \item \textbf{Sur-\'echantillonnage} (over)~: tirage avec remise
    dans la classe minoritaire. Risque de sur-ajustement (doublons).
  \item \textbf{Sous-\'echantillonnage} (under)~: tirage sans remise
    dans la classe majoritaire. Perte d'information.
  \item \textbf{Combinaison} (both)~: sur- et sous-\'echantillonnage
    simultan\'es.
\end{itemize}

\begin{table}[h]
\centering
\caption{Impact de l'\'equilibrage sur l'AUC (hacide, arbre rpart)}
\begin{tabular}{lcc}
\toprule
M\'ethode & AUC & Commentaire \\
\midrule
Aucun (arbre brut) & 0,600 & Mauvais, classe rare ignor\'ee \\
Oversampling & 0,798 & Am\'elioration significative \\
Undersampling & 0,749 & Moindre que l'over \\
Both & 0,798 & \'Equivalent \`a l'over \\
\textbf{ROSE} & \textbf{0,989} & \textbf{Excellente g\'en\'eralisation} \\
\bottomrule
\end{tabular}
\end{table}

% ===================================================================
\section{La m\'ethode ROSE}
% ===================================================================

\subsection{Algorithme}

ROSE g\'en\`ere un nouvel exemple artificiel $(x^*, y^*)$ en trois
\'etapes~:

\begin{enumerate}
  \item Choisir $y^* = Y_j$ avec probabilit\'e $\pi_j$.
  \item S\'electionner $(x_i, y_i)$ tel que $y_i = y^*$,
    avec probabilit\'e $1/n_j$.
  \item G\'en\'erer $x^* \sim K_{H_j}(\cdot, x_i)$, o\`u $K_{H_j}$
    est un noyau centr\'e en $x_i$ de matrice de lissage $H_j$.
\end{enumerate}

\subsection{Propri\'et\'es}

\begin{itemize}
  \item \'Equivalent \`a tirer de l'estimateur \`a noyau de
    $f(x|Y_j)$ (densit\'e conditionnelle).
  \item Quand $H_j \to 0$, ROSE se r\'eduit \`a une combinaison
    over/under-sampling.
  - Pas de doublons (contrairement au sur-\'echantillonnage simple).
  \item Les param\`etres de lissage $H_0$, $H_1$ sont choisis de
    fa\c con optimale sous hypoth\`ese de normalit\'e.
  \item Ajustables via \texttt{hmult.majo} et \texttt{hmult.mino}.
\end{itemize}

\subsection{ROSE pour l'\'evaluation}

ROSE peut \^etre utilis\'e pour estimer la pr\'ecision d'un
classifieur selon trois sch\'emas~:

\begin{itemize}
  \item \textbf{Holdout}~: entra\^inement sur $\mathcal{T}^*_m$
    (ROSE), test sur $\mathcal{T}_n$ (original).
  \item \textbf{Bootstrap}~: $B$ \'echantillons ROSE, test sur
    l'original, distribution bootstrap de la m\'etrique.
  \item \textbf{Leave-$K$-out CV}~: $Q = n/K$ groupes, entra\^inement
    ROSE sur $\mathcal{T}_n \setminus \mathcal{T}^K_i$, test sur
    $\mathcal{T}^K_i$.
\end{itemize}

\subsection{Analyse critique}

\textbf{Forces~:}
\begin{itemize}
  \item Approche unifi\'ee (estimation + \'evaluation).
  \item Pas de doublons, meilleure g\'en\'eralisation.
  \item Flexible (tout classifieur R standard ou sur mesure).
  \item Interface bien con\c cue, documentation claire.
\end{itemize}

\textbf{Faiblesses~:}
\begin{itemize}
  \item Param\`etres de lissage $H_j$ choisis sous hypoth\`ese de
    normalit\'e (potentiellement sous-optimaux pour d'autres
    distributions).
  \item Co\^ut calculatoire pour bootstrap $B=50$ (50 entra\^\i nements
    + pr\'edictions).
  \item Pas de parall\'elisation native.
  \item M\'ethode limit\'ee \`a la classification binaire.
  \item Holdout ROSE tend \`a surestimer l'AUC (biais optimiste),
    mais avec un MSE plus faible qu'un holdout standard.
\end{itemize}

% ===================================================================
\section{Illustrations num\'eriques}
% ===================================================================

\subsection{Donn\'ees hacide}

Les donn\'ees \texttt{hacide} sont fournies avec le package~:
\begin{itemize}
  \item \texttt{hacide.train}~: $n=1000$, 2\% de positifs (20/980).
  \item \texttt{hacide.test}~: $n=250$.
  \item 2 variables explicatives $x_1$, $x_2$.
  \item La classe rare forme un demi-cercle bruit\'e au sein de la
    classe majoritaire (ellipse normale).
\end{itemize}

\subsection{R\'esultats comparatifs}

\begin{figure}[h]
\centering
\fbox{\parbox{0.85\textwidth}{\centering
\textbf{AUC sur l'\'echantillon de test hacide.test} \\
\vspace{0.3cm}
\begin{tabular}{lcc}
\toprule
Strat\'egie & AUC & \'Ecart / base \\
\midrule
Arbre brut (d\'es\'equilibr\'e) & 0,600 & --- \\
Arbre + oversampling & 0,798 & $+0,198$ \\
Arbre + undersampling & 0,749 & $+0,149$ \\
Arbre + both & 0,798 & $+0,198$ \\
\textbf{Arbre + ROSE} & \textbf{0,989} & $\mathbf{+0,389}$ \\
\midrule
ROSE holdout (AUC estim\'ee) & 0,985 & --- \\
ROSE bootstrap (m\'ediane) & 0,984 & --- \\
ROSE L5OCV & 0,973 & --- \\
\bottomrule
\end{tabular}
}}
\caption{Performances compar\'ees des strat\'egies d'\'equilibrage
sur les donn\'ees hacide (Lunardon et al., 2014)}
\end{figure}

\newpage
% ===================================================================
\section{Extraits de code R}
% ===================================================================

\subsection{Installation et chargement}

\begin{lstlisting}
install.packages("ROSE")
library(ROSE)
data(hacide)
str(hacide.train)
table(hacide.train$cls)  # 0: 980, 1: 20
\end{lstlisting}

\subsection{Arbre na\"if (ignorant le d\'es\'equilibre)}

\begin{lstlisting}
library(rpart)
treeimb <- rpart(cls ~ ., data = hacide.train)
pred.treeimb <- predict(treeimb, newdata = hacide.test)

# Metriques
accuracy.meas(hacide.test$cls, pred.treeimb[,2])
# precision: 1.000, recall: 0.200, F: 0.167

roc.curve(hacide.test$cls, pred.treeimb[,2], plotit = FALSE)
# AUC: 0.600
\end{lstlisting}

\subsection{Sur-\'echantillonnage (ovun.sample)}

\begin{lstlisting}
# Oversampling
data.bal.ov <- ovun.sample(cls ~ ., data = hacide.train,
                            method = "over", p = 0.5, seed = 1)$data
table(data.bal.ov$cls)

# Undersampling
data.bal.un <- ovun.sample(cls ~ ., data = hacide.train,
                            method = "under", p = 0.5, seed = 1)$data
table(data.bal.un$cls)

# Combinaison
data.bal.ou <- ovun.sample(cls ~ ., data = hacide.train,
                            method = "both", N = 1000, p = 0.5,
                            seed = 1)$data
table(data.bal.ou$cls)
\end{lstlisting}

\subsection{M\'ethode ROSE}

\begin{lstlisting}
# Generation de donnees equilibrees
data.rose <- ROSE(cls ~ ., data = hacide.train, seed = 1)$data
table(data.rose$cls)

# Ajustement et prediction
tree.rose <- rpart(cls ~ ., data = data.rose)
pred.tree.rose <- predict(tree.rose, newdata = hacide.test)
roc.curve(hacide.test$cls, pred.tree.rose[,2])
# AUC: 0.989

# Ajustement des parametres de lissage
data.rose.h <- ROSE(cls ~ ., data = hacide.train, seed = 1,
                    hmult.majo = 0.25, hmult.mino = 0.5)$data
\end{lstlisting}

\subsection{ROSE pour l'\'evaluation}

\begin{lstlisting}
# Holdout ROSE
ROSE.holdout <- ROSE.eval(cls ~ ., data = hacide.train,
    learner = rpart, method.assess = "holdout",
    extr.pred = function(obj) obj[,2], seed = 1)
# AUC: 0.985

# Bootstrap ROSE (B=50)
ROSE.boot <- ROSE.eval(cls ~ ., data = hacide.train,
    learner = rpart, method.assess = "BOOT", B = 50,
    extr.pred = function(obj) obj[,2], seed = 1)
summary(ROSE.boot)
# Mediane AUC: 0.984

# L5OCV ROSE
ROSE.l5ocv <- ROSE.eval(cls ~ ., data = hacide.train,
    learner = rpart, method.assess = "LKOCV", K = 5,
    extr.pred = function(obj) obj[,2], seed = 1)
# L5OCV AUC: 0.973
\end{lstlisting}

\subsection{Utilisation d'un classifieur non-standard (k-NN)}

\begin{lstlisting}
library(class)
knn.wrap <- function(data, newdata, ...) {
  knn(train = data[,-1], test = newdata, cl = data[,1], ...)
}

rose.holdout.knn <- ROSE.eval(cls ~ ., data = hacide.train,
    learner = knn.wrap, control.learner = list(k = 2, prob = TRUE),
    method.assess = "holdout", seed = 1)
\end{lstlisting}

% ===================================================================
\section{Conclusion}
% ===================================================================

Le package \textbf{ROSE} est un outil complet et bien con\c cu pour
la classification binaire sur donn\'ees d\'es\'equilibr\'ees. Ses
points forts~:
\begin{itemize}
  \item M\'ethode unifi\'ee couvrant estimation et \'evaluation.
  \item G\'en\'eration de donn\'ees artificielles sans doublons
    (bootstrap liss\'e).
  \item Interface souple acceptant tout classifieur R.
  \item Plusieurs m\'ethodes d'\'evaluation (holdout, bootstrap, CV).
\end{itemize}

Ses limites~:
\begin{itemize}
  \item Classification binaire uniquement.
  \item Param\'etrisation du noyau sous hypoth\`ese de normalit\'e.
  \item Co\^ut calculatoire non n\'egligeable (surtout bootstrap).
  \item L\'eg\`ere surestimation de l'AUC en holdout.
\end{itemize}

\textbf{Note~:} Depuis 2014, d'autres packages ont \'emerg\'e
(\texttt{themis} pour le tidyverse, \texttt{smotefamily},
\texttt{UBL}), mais ROSE reste une solution de r\'ef\'erence,
p\'edagogique et robuste.

% ===================================================================
\newpage
\section*{Questions cl\'es (QCM)}
\addcontentsline{toc}{section}{Questions cl\'es (QCM)}

\begin{enumerate}

  \item Quel est le principal probl\`eme pos\'e par les donn\'ees
    d\'es\'equilibr\'ees en classification~?
    \begin{enumerate}
      \item Le mod\`ele converge trop vite
      \item La classe rare est ignor\'ee par le classifieur
      \item Les variables explicatives sont corr\'el\'ees
      \item L'AUC est toujours de 0,5
    \end{enumerate}
    \textbf{R\'eponse~: b.}

  \item Dans l'exemple hacide, quel est le pourcentage d'exemples
    positifs dans l'\'echantillon d'apprentissage~?
    \begin{enumerate}
      \item 10\%
      \item 5\%
      \item 2\%
      \item 20\%
    \end{enumerate}
    \textbf{R\'eponse~: c (20/1000 = 2\%).}

  \item Quelle est la principale diff\'erence entre ROSE et le
    sur-\'echantillonnage simple~?
    \begin{enumerate}
      \item ROSE est plus lent
      \item ROSE g\'en\`ere des donn\'ees synth\'etiques sans
        doublons
      \item ROSE n\'ecessite la library caret
      \item ROSE ne fonctionne qu'avec rpart
    \end{enumerate}
    \textbf{R\'eponse~: b.}

  \item Dans l'algorithme ROSE, \`a quoi correspond la matrice
    $H_j$~?
    \begin{enumerate}
      \item \`A la matrice de covariance des donn\'ees
      \item \`A la matrice de lissage du noyau
      \item \`A la matrice des poids du classifieur
      \item \`A la matrice de confusion
    \end{enumerate}
    \textbf{R\'eponse~: b.}

  \item Quelle AUC obtient-on avec un arbre de d\'ecision
    (\texttt{rpart}) entra\^\i n\'e sur des donn\'ees \'equilibr\'ees
    par ROSE (donn\'ees hacide)~?
    \begin{enumerate}
      \item 0,600
      \item 0,798
      \item 0,989
      \item 0,500
    \end{enumerate}
    \textbf{R\'eponse~: c.}

  \item Quelle fonction du package ROSE permet d'obtenir une
    estimation de l'AUC par validation crois\'ee~?
    \begin{enumerate}
      \item \texttt{ROSE()}
      \item \texttt{ovun.sample()}
      \item \texttt{ROC.eval()}
      \item \texttt{ROSE.eval()}
    \end{enumerate}
    \textbf{R\'eponse~: d.}

  \item Quelle est la m\'etrique d'\'evaluation recommand\'ee par
    les auteurs pour \'eviter le choix arbitraire d'un
    seuil~?
    \begin{enumerate}
      \item Le $F$-mesure
      \item La pr\'ecision
      \item Le rappel
      \item L'AUC (aire sous la courbe ROC)
    \end{enumerate}
    \textbf{R\'eponse~: d.}

  \item Quel argument permet de r\'eduire la largeur du noyau dans
    \texttt{ROSE()}~?
    \begin{enumerate}
      \item \texttt{hmult.majo} et \texttt{hmult.mino}
      \item \texttt{kernel.width}
      \item \texttt{smoothing.factor}
      \item \texttt{bandwidth}
    \end{enumerate}
    \textbf{R\'eponse~: a.}

  \item Que se passe-t-il quand $H_j \to 0$ dans la m\'ethode
    ROSE~?
    \begin{enumerate}
      \item La m\'ethode diverge
      \item ROSE se r\'eduit \`a over/under-sampling standard
      \item Le noyau devient gaussien
      \item AUC devient 1
    \end{enumerate}
    \textbf{R\'eponse~: b.}

  \item Selon l'article, comment peut-on adapter un classifieur
    non-standard (ex.~\texttt{knn}) pour l'utiliser dans
    \texttt{ROSE.eval()}~?
    \begin{enumerate}
      \item Ce n'est pas possible
      \item En \'ecrivant une fonction wrapper avec arguments
        \texttt{data} et \texttt{newdata}
      \item En modifiant le code source de ROSE
      \item En utilisant \texttt{caret::train} \`a la place
    \end{enumerate}
    \textbf{R\'eponse~: b.}

\end{enumerate}

% ===================================================================
\begin{thebibliography}{5}
% ===================================================================

\bibitem{Lunardon2014}
  N.~Lunardon, G.~Menardi, N.~Torelli.
  \newblock \emph{ROSE: A Package for Binary Imbalanced Learning}.
  \newblock The R Journal, 6(1):79--89, 2014.

\bibitem{Menardi2014}
  G.~Menardi, N.~Torelli.
  \newblock \emph{Training and assessing classification rules with
    imbalanced data}.
  \newblock Data Mining and Knowledge Discovery, 28(1):92--122, 2014.

\bibitem{Chawla2002}
  N.~V.~Chawla, K.~W.~Bowyer, L.~O.~Hall, W.~P.~Kegelmeyer.
  \newblock \emph{SMOTE: Synthetic Minority Over-sampling Technique}.
  \newblock J.~Artif.~Intell.~Res., 16:321--357, 2002.

\bibitem{Torgo2010}
  L.~Torgo.
  \newblock \emph{Data Mining with R, Learning with Case Studies}.
  \newblock Chapman \& Hall/CRC, 2010.

\end{thebibliography}

\end{document}
